This chapter introduces the research context, motivation, and scope of the FreelanceChain project. Section \ref{section:1.1} identifies the structural limitations of existing freelance platforms that motivate a decentralized solution. Section \ref{section:1.2} defines the objectives and scope of this thesis. Section \ref{section:1.3} outlines the proposed technical approach, and Section \ref{section:1.4} describes the organization of subsequent chapters.

\section{Motivation}
\label{section:1.1}

The global gig economy has grown at an unprecedented rate over the past decade. By 2023, the online freelance market was estimated to be worth approximately \$455 billion, with over 1.5 billion independent workers participating worldwide \cite{statista2024freelance}. Digital platforms that match freelancers with clients have been central to this growth, with Upwork reporting 868,000 active clients and Fiverr processing millions of transactions annually \cite{upwork2023, fiverr2023}. However, despite their scale, these platforms share a set of structural deficiencies that impose significant costs on participants.

\textbf{The intermediary problem.} Every interaction on a centralized freelance platform is mediated by a private corporation acting as a trusted third party. This corporation holds escrow funds, enforces contracts, resolves disputes, and sets platform rules. While this arrangement provides convenience, it concentrates enormous power in a single actor. Platforms routinely charge service fees between 5\% and 20\% of the transaction value---fees that directly reduce freelancer income and inflate client costs. In 2022 alone, Upwork collected over \$600 million in revenue, the majority of which derived from these intermediary fees.

\textbf{Payment insecurity and fraud.} Freelancers operating on centralized platforms face systemic payment risk. Platform terms-of-service allow accounts to be suspended without prior notice, instantly freezing any held funds. Clients frequently exploit dispute mechanisms to avoid payment for completed work, knowing that the platform's adjudication process favors the paying party. A 2023 survey found that over 34\% of freelancers reported at least one non-payment incident in the preceding twelve months.

\textbf{Opacity of reputation.} Rating systems on centralized platforms are maintained in proprietary databases controlled by the platform operator. A freelancer who builds a five-star reputation on Upwork cannot transfer that reputation to Fiverr or any other platform. Furthermore, these systems are subject to manipulation: fake reviews, suppressed negative feedback, and algorithmic biases distort the signal they are meant to convey. Clients cannot independently verify the authenticity of a freelancer's stated track record.

\textbf{Identity and Sybil attacks.} The pseudonymous nature of online accounts makes it trivial for malicious actors to create multiple fake identities (a so-called Sybil attack) to game rating systems, submit fraudulent bids, or conduct payment fraud. Centralized platforms address this through proprietary identity verification processes, but the data collected is stored in centralized databases that are prime targets for breaches.

\textbf{The blockchain opportunity.} The advent of public smart-contract blockchains offers a fundamentally different architecture for multi-party trust. Smart contracts, first articulated by Nick Szabo \cite{szabo1994smart} and implemented at scale by Ethereum \cite{buterin2014ethereum}, allow contractual obligations to be encoded in self-executing code that runs without a trusted administrator. Funds locked in a smart contract escrow cannot be seized by a platform operator; rules encoded in bytecode cannot be changed retroactively; and all transactions are publicly auditable on an immutable ledger.

The Solana blockchain \cite{yakovenko2018solana} further reduces the economic barrier to blockchain adoption by providing transaction throughput exceeding 50,000 \gls{tps} at a cost below \$0.001 per transaction---compared to Ethereum's occasional fees exceeding \$50 per transaction. This performance profile makes Solana the most viable public blockchain for consumer-facing applications that require frequent, low-value transactions.

These observations together motivate the development of a fully decentralized freelance marketplace that eliminates the intermediary, secures payments in tamper-proof escrow, records reputation on-chain as non-transferable tokens, and provides identity verification without centralizing sensitive biometric data.

\section{Objectives and Scope of the Graduation Thesis}
\label{section:1.2}

Existing research and deployed systems in the space of decentralized freelancing are limited in scope. Brondeau et al. \cite{delmolino2016step} demonstrated the feasibility of simple escrow smart contracts but did not address reputation, identity, or governance. Platforms such as Ethlance and Colony provide basic on-chain job posting and payment but lack biometric identity verification, AI-assisted matching, milestone-based fund release, and structured governance. No existing system addresses the full lifecycle of a professional freelance engagement in a decentralized manner.

The present limitations can be summarized as follows:
\begin{enumerate}
    \item No comprehensive smart contract suite covers the complete job lifecycle from creation to reputation recording.
    \item No existing decentralized platform provides privacy-preserving biometric identity verification with on-chain attestation.
    \item No platform combines Soulbound Token reputation, zero-knowledge credential proofs, AI risk assessment, and DAO governance in a unified system.
    \item High transaction costs on Ethereum-based platforms make micro-payment freelancing economically infeasible.
\end{enumerate}

This thesis addresses these limitations by designing, implementing, and deploying FreelanceChain --- a comprehensive decentralized freelance marketplace on Solana. The specific objectives of this work are:

\begin{enumerate}
    \item Design and implement an Anchor smart contract suite covering job management, dual-track escrow (lump sum and milestone-based), on-chain reputation via \gls{sbt}, biometric \gls{kyc}, \gls{dao} governance, social recovery, zero-knowledge credential verification, and AI oracle consensus.
    \item Build a production-quality Next.js frontend that provides a seamless user experience for both clients and freelancers, integrating Solana wallet adapters and supporting guest browsing.
    \item Integrate browser-side face recognition to enable privacy-preserving identity verification without uploading biometric data to any server.
    \item Demonstrate the economic and technical viability of the system through deployment on the Solana Devnet and systematic testing.
\end{enumerate}

The scope of this thesis is bounded to the Solana Devnet deployment. Mainnet launch, regulatory compliance analysis, and large-scale user studies are identified as future work.

\section{Tentative Solution}
\label{section:1.3}

FreelanceChain adopts a blockchain-first architecture in which all trust-critical operations --- escrow, reputation, identity attestation, and governance --- execute on-chain via Anchor smart contracts. The following design principles guide the solution:

\textbf{Smart contract-enforced escrow.} Funds are locked in Program Derived Address (\gls{pda}) accounts controlled exclusively by contract logic. Neither the platform developer nor any third party can access locked funds. Funds are released only when the conditions encoded in the contract are satisfied: client approval, milestone completion, or arbitrator ruling.

\textbf{Soulbound Token reputation.} Following the model proposed by Buterin, Weyl, and Ohlhaver \cite{weyl2022decentralized}, reputation ratings are minted as non-transferable \gls{sbt} tokens attached to the wallet that earned them. Each token carries a rating (1--5), a comment, a SHA-256 verification hash, and immutable metadata. The non-transferability property prevents the reputation market from being gamed by buying and selling established accounts.

\textbf{Browser-side biometric KYC.} Face recognition is performed entirely in the user's browser using the \texttt{@vladmandic/face-api} library with TinyFaceDetector and FaceRecognition models. The user photographs their identity document and captures a webcam selfie; the library computes the Euclidean distance between the two face descriptors. If this distance is below the threshold of 0.50, the user is considered verified. Only the match result and distance (stored as basis points) are written to the blockchain --- no image or biometric feature vector is transmitted to any server.

\textbf{AI-assisted risk assessment.} A Next.js serverless API route invokes the \texttt{llama-3.3-70b-versatile} \gls{llm} via the Groq API to analyze CV text submitted alongside a bid. The model produces a structured four-score assessment (skill match, CV authenticity, budget fit, and timeline realism), a composite risk classification (LOW / MEDIUM / HIGH), and supports follow-up questions through a conversational AI advisor. This assists clients in evaluating candidates at scale without replacing human judgment.

\textbf{DAO governance.} Platform parameters (fee rates, dispute thresholds, governance quorum) are controlled by on-chain \gls{dao} proposals. Token holders create proposals by staking governance tokens and vote with reputation-weighted ballots. This ensures that platform evolution is subject to community consensus rather than unilateral operator decisions.

The choice of Solana is justified by its proof-of-history consensus mechanism, which achieves sub-second finality at a fraction of the cost of proof-of-work or proof-of-stake alternatives on Ethereum. The Anchor framework provides a strongly-typed Rust abstraction that significantly reduces smart contract vulnerability surface compared to raw Solana program development.

\section{Thesis Organization}
\label{section:1.4}

The remainder of this thesis is organized as follows.

Chapter 2 surveys the current state of freelance platforms, analyzes the deficiencies of centralized alternatives, identifies the actors and use cases of FreelanceChain, and specifies both functional and non-functional requirements. The chapter provides the stakeholder perspective and use case specifications that drive all subsequent design decisions.

Chapter 3 presents the methodology and technologies employed in FreelanceChain. It covers the Solana blockchain and its proof-of-history consensus mechanism, the Anchor smart contract framework, the Rust programming language, the Next.js and React frontend stack, face recognition with face-api.js, \gls{llm}-powered risk assessment, and the rationale for choosing each technology over its principal alternatives.

Chapter 4 details the system architecture, the design of the smart contract modules and their \gls{pda} account schemas, the frontend component hierarchy, the database and storage design using \gls{ipfs} and browser localStorage, the tools and libraries used, and the testing methodology applied to the critical transaction paths.

Chapter 5 presents the principal technical contributions of this thesis: the dual-track escrow model, the Soulbound Token reputation system, the browser-side biometric KYC pipeline, the AI oracle integration, and the DAO governance mechanism. Each contribution is analyzed in terms of the problem it addresses, the solution implemented, and the results achieved.

Chapter 6 concludes the thesis by comparing FreelanceChain against existing systems, summarizing achieved and unachieved objectives, and identifying future research and development directions.
